\documentclass[11pt,a4paper]{article}
\usepackage[utf8]{inputenc}
\usepackage[T1]{fontenc}
\usepackage{amsmath, amssymb, amsthm}
\usepackage{geometry}
\geometry{margin=1in}
\usepackage{fancyhdr}
\usepackage{listings}
\usepackage{newunicodechar}
\usepackage{hyperref}

\pagestyle{fancy}
\fancyhf{}
\rhead{Conjecture d'Erd\H{o}s-Straus}
\cfoot{\thepage}

\title{Formulation Rigoureuse et Strat\'egie de Preuve pour la Conjecture d'Erd\H{o}s-Straus}
\author{Charles EDOU NZE, chercheur ind\'ependant}
\date{\today}

\newtheorem{theorem}{Th\'eor\`eme}[section]
\newtheorem{lemma}[theorem]{Lemme}
\newtheorem{definition}[theorem]{D\'efinition}
\newtheorem{proposition}[theorem]{Proposition}

\newunicodechar{ℕ}{\ensuremath{\mathbb{N}}}
\newunicodechar{∀}{\ensuremath{\forall}}
\newunicodechar{∃}{\ensuremath{\exists}}
\newunicodechar{∧}{\ensuremath{\wedge}}
\newunicodechar{∨}{\ensuremath{\vee}}
\newunicodechar{→}{\ensuremath{\rightarrow}}

\lstset{literate=
  {é}{{\'e}}1
  {è}{{\`e}}1
  {ê}{{\^e}}1
  {à}{{\`a}}1
}

\begin{document}

\maketitle

\begin{abstract}
La conjecture d'Erd\H{o}s-Straus postule que pour tout entier $n > 1$, la fraction $4/n$ peut \^etre exprim\'ee comme la somme de trois fractions unitaires positives. Ce document \'etablit une fondation axiomatique rigoureuse, passe en revue la litt\'erature pertinente concernant les identit\'es modulaires, et fournit une d\'erivation math\'ematique explicite pour deux cas modulaires sp\'ecifiques ($n \equiv 2 \pmod 3$ et $n \equiv 3 \pmod 4$). En outre, une architecture structurelle pour l'autoformalisation dans l'assistant de preuve Lean 4 est pr\'esent\'ee.

\vspace{2cm}
\begin{flushright}
\textit{Charles EDOU NZE, chercheur ind\'ependant}
\end{flushright}
\end{abstract}

\newpage

\section{D\'efinitions Axiomatiques}

\begin{definition}
Soit $\mathbb{N}_{>1} = \{n \in \mathbb{N} \mid n \geq 2\}$. La conjecture d'Erd\H{o}s-Straus postule que pour tout $n \in \mathbb{N}_{>1}$, il existe $x, y, z \in \mathbb{N}_{>0}$ tels que :
\begin{equation}
\frac{4}{n} = \frac{1}{x} + \frac{1}{y} + \frac{1}{z}
\end{equation}
\end{definition}

\begin{definition}
En multipliant l'\'equation susmentionn\'ee par $n x y z$, nous obtenons la formulation diophantienne polynomiale \'equivalente :
\begin{equation}
4 x y z = n (x y + x z + y z)
\end{equation}
\end{definition}

\section{Litt\'erature Contextuelle et Travaux Connexes}

La conjecture d'Erd\H{o}s-Straus, formul\'ee en 1948 par Paul Erd\H{o}s et Ernst G. Straus, cherche une expansion de $4/n$ en exactement trois fractions \'egyptiennes positives.
La litt\'erature pertinente explore de mani\`ere approfondie les identit\'es modulaires. Par exemple, Mordell (1967) a d\'emontr\'e qu'une identit\'e fournissant des solutions pour tous les $n \equiv r \pmod p$ existe si et seulement si $r$ n'est pas un r\'esidu quadratique modulo $p$.
Des recherches informatiques (par exemple, par Salez, 2014) ont v\'erifi\'e la conjecture jusqu'\`a $N = 10^{17}$.
Elsholtz et Tao (2013) ont born\'e le nombre moyen de solutions de fa\c{c}on polylogarithmique en utilisant les th\'eor\`emes de Bombieri-Vinogradov et de Brun-Titchmarsh.

Le principe de Hasse sugg\`ere que l'\'etude des \'equations diophantiennes modulo $p$ pourrait permettre d'assembler des solutions enti\`eres. Cependant, un syst\`eme complet de congruences couvrantes est av\'er\'e impossible car 1 est un r\'esidu quadratique modulo n'importe quel $n > 1$.

\section{Strat\'egie de Preuve et Lemmes}

Le probl\`eme peut \^etre r\'eduit \`a la consid\'eration des seules valeurs premi\`eres de $n$. Si $n = m \cdot p$ o\`u $p$ est premier, une solution pour $p$ donne une solution pour $n$ en redimensionnant les d\'enominateurs par $m$.

\begin{lemma}[Cas Modulaire $n \equiv 2 \pmod 3$]
Pour tout entier $n > 1$ tel que $n \equiv 2 \pmod 3$, il existe une solution $(x,y,z) \in \mathbb{N}_{>0}^3$ satisfaisant $\frac{4}{n} = \frac{1}{x} + \frac{1}{y} + \frac{1}{z}$.
\end{lemma}

\begin{lemma}[Cas Modulaire $n \equiv 3 \pmod 4$]
Pour tout entier $n > 1$ tel que $n \equiv 3 \pmod 4$, il existe une solution $(x,y,z) \in \mathbb{N}_{>0}^3$ satisfaisant $\frac{4}{n} = \frac{1}{x} + \frac{1}{y} + \frac{1}{z}$.
\end{lemma}

\section{Preuve Math\'ematique (Z\'ero Ellipse)}

Nous proc\'edons \`a la preuve explicite des lemmes.

\begin{proof}[Preuve du Lemme 3.1]
Soit $n \in \mathbb{N}_{>1}$ tel que $n \equiv 2 \pmod 3$.
Par d\'efinition de la congruence modulo 3, il existe un entier $k \geq 0$ tel que $n = 3k + 2$.
Nous d\'efinissons $x = n$. Puisque $n \geq 2$, $x \in \mathbb{N}_{>0}$.
Nous d\'efinissons $y = \frac{n+1}{3}$.
Nous devons v\'erifier que $y$ est un entier.
Substituons $n = 3k + 2$ dans l'expression pour $y$ :
\begin{align*}
y &= \frac{(3k + 2) + 1}{3} \\
  &= \frac{3k + 3}{3} \\
  &= k + 1
\end{align*}
Puisque $k \geq 0$, $k+1$ est un entier et $y \in \mathbb{N}_{>0}$.
Nous d\'efinissons $z = \frac{n(n+1)}{3}$.
Nous devons v\'erifier que $z$ est un entier.
Substituons $n = 3k + 2$ :
\begin{align*}
z &= \frac{(3k+2)(3k+3)}{3} \\
  &= (3k+2)(k+1)
\end{align*}
Puisque $k \geq 0$, $3k+2 \geq 2$ et $k+1 \geq 1$. Les deux sont des entiers, donc leur produit est un entier, et $z \in \mathbb{N}_{>0}$.

Nous calculons maintenant la somme $\frac{1}{x} + \frac{1}{y} + \frac{1}{z}$ :
\begin{align*}
\frac{1}{x} + \frac{1}{y} + \frac{1}{z} &= \frac{1}{n} + \frac{1}{\frac{n+1}{3}} + \frac{1}{\frac{n(n+1)}{3}} \\
&= \frac{1}{n} + \frac{3}{n+1} + \frac{3}{n(n+1)}
\end{align*}
Pour additionner ces fractions, nous trouvons un d\'enominateur commun, qui est $n(n+1)$.
\begin{align*}
\frac{1}{n} &= \frac{n+1}{n(n+1)} \\
\frac{3}{n+1} &= \frac{3n}{n(n+1)} \\
\frac{3}{n(n+1)} &= \frac{3}{n(n+1)}
\end{align*}
En sommant les num\'erateurs :
\begin{align*}
\frac{n+1 + 3n + 3}{n(n+1)} &= \frac{4n + 4}{n(n+1)} \\
&= \frac{4(n+1)}{n(n+1)}
\end{align*}
Nous divisons le num\'erateur et le d\'enominateur par $(n+1)$. Puisque $n \geq 2$, $n+1 \neq 0$.
\begin{align*}
\frac{4(n+1)}{n(n+1)} = \frac{4}{n}
\end{align*}
Ainsi, $\frac{1}{x} + \frac{1}{y} + \frac{1}{z} = \frac{4}{n}$.
Ceci ach\`eve la preuve.
\end{proof}

\begin{proof}[Preuve du Lemme 3.2]
Soit $n \in \mathbb{N}_{>1}$ tel que $n \equiv 3 \pmod 4$.
Par d\'efinition, il existe un entier $k \geq 0$ tel que $n = 4k + 3$.
Nous d\'efinissons $x = \frac{n+1}{4}$.
Substituons $n = 4k + 3$ :
\begin{align*}
x &= \frac{(4k+3)+1}{4} \\
  &= \frac{4k+4}{4} \\
  &= k+1
\end{align*}
Puisque $k \geq 0$, $x \in \mathbb{N}_{>0}$.
Nous d\'efinissons $y = \frac{n(n+1)}{2}$.
Substituons $n = 4k + 3$ :
\begin{align*}
y &= \frac{(4k+3)(4k+4)}{2} \\
  &= \frac{(4k+3)4(k+1)}{2} \\
  &= 2(4k+3)(k+1)
\end{align*}
Puisque $k \geq 0$, $4k+3 \geq 3$ et $k+1 \geq 1$. Les deux sont des entiers, donc leur produit multipli\'e par 2 est un entier, et $y \in \mathbb{N}_{>0}$.
Nous d\'efinissons $z = \frac{n(n+1)}{2}$.
Puisque cette expression est identique \`a $y$, $z \in \mathbb{N}_{>0}$.

Nous calculons maintenant la somme $\frac{1}{x} + \frac{1}{y} + \frac{1}{z}$ :
\begin{align*}
\frac{1}{x} + \frac{1}{y} + \frac{1}{z} &= \frac{1}{\frac{n+1}{4}} + \frac{1}{\frac{n(n+1)}{2}} + \frac{1}{\frac{n(n+1)}{2}} \\
&= \frac{4}{n+1} + \frac{2}{n(n+1)} + \frac{2}{n(n+1)}
\end{align*}
En sommant les deux derni\`eres fractions :
\begin{align*}
\frac{2}{n(n+1)} + \frac{2}{n(n+1)} &= \frac{4}{n(n+1)}
\end{align*}
Maintenant, sommons avec la premi\`ere fraction :
\begin{align*}
\frac{4}{n+1} + \frac{4}{n(n+1)}
\end{align*}
Trouvons le d\'enominateur commun $n(n+1)$ :
\begin{align*}
\frac{4}{n+1} &= \frac{4n}{n(n+1)}
\end{align*}
En sommant les num\'erateurs :
\begin{align*}
\frac{4n + 4}{n(n+1)} &= \frac{4(n+1)}{n(n+1)}
\end{align*}
Divisons le num\'erateur et le d\'enominateur par $(n+1)$. Puisque $n \geq 2$, $n+1 \neq 0$ :
\begin{align*}
\frac{4(n+1)}{n(n+1)} = \frac{4}{n}
\end{align*}
Ainsi, $\frac{1}{x} + \frac{1}{y} + \frac{1}{z} = \frac{4}{n}$.
Ceci ach\`eve la preuve.
\end{proof}

\section{Architecture pour l'Autoformalisation en Lean 4}

Pour faciliter la v\'erification automatis\'ee par un syst\`eme formel tel que Lean 4, les preuves peuvent \^etre codifi\'ees comme suit.

\begin{lstlisting}[basicstyle=\ttfamily\small, breaklines=true]
import Mathlib.Data.Nat.Basic
import Mathlib.Data.Nat.Parity

/--
La Conjecture d'Erdos-Straus
-/
def erdos_straus_conjecture : Prop :=
  ∀ n : ℕ, n > 1 → ∃ x y z : ℕ, x > 0 ∧ y > 0 ∧ z > 0 ∧
    4 * x * y * z = n * (x * y + x * z + y * z)

/--
Lemme 1 : Solution pour n = 2 mod 3
-/
theorem erdos_straus_mod3_2 (n : ℕ) (h1 : n > 1) (h2 : n % 3 = 2) :
  ∃ x y z : ℕ, x > 0 ∧ y > 0 ∧ z > 0 ∧
    4 * x * y * z = n * (x * y + x * z + y * z) := by
  admit

/--
Lemme 2 : Solution pour n = 3 mod 4
-/
theorem erdos_straus_mod4_3 (n : ℕ) (h1 : n > 1) (h2 : n % 4 = 3) :
  ∃ x y z : ℕ, x > 0 ∧ y > 0 ∧ z > 0 ∧
    4 * x * y * z = n * (x * y + x * z + y * z) := by
  admit
\end{lstlisting}

\end{document}
